\chapter{Conclusions}
\label{ch:conclusions}

\section{Summary}
\label{sec:conclusion-summary}

This study assessed the state of the practice for robotics motion planning and
simulation (MPS) software by adapting the Domain~X methodology
(Chapter~\ref{ch:methodology}) to a new domain. From 295 consolidated candidate
names, 28 packages were selected and measured in May 2026 against the Domain~X
template, spanning full simulators, physics and dynamics engines, planning
libraries, manipulation frameworks, and middleware. A commonality analysis
(Chapter~\ref{ch:selected-software-commonality}) described the selected packages
as a software family, and the findings were compared with prior studies of the
state of the practice for medical imaging and Lattice Boltzmann method software.
A developer-interview component was designed and received MREB ethics clearance
(project~\#8215, 2026-06-22); recruitment is underway, but the interviews have
not yet been conducted, so their evidence is not reported here.

\section{Answers to Research Questions}
\label{sec:answers-research-questions}

\textbf{RQ1 (candidate packages).} Filtering the
295 candidates by open-source availability, repository-based measurability,
activity status, and scope coherence (Chapter~\ref{ch:methodology},
Appendix~\ref{appendix:full-candidate-software-list}) yielded the 28 measured packages; notable
exclusions include MATLAB/Simulink, Unity, NVIDIA Isaac Sim, and RaiSim (closed
core components), MRDS and USARSim (inactive), and GPMP2 and CHOMP
(single-algorithm scope).

\textbf{RQ2 (quality ranking).} Across the nine quality categories
(Chapter~\ref{ch:measurement-results}), the unweighted quality mean is led by
Gazebo (10.0), Drake (8.8), and a tied group of MoveIt!, OMPL, Pinocchio, and
ROS~2 (8.7); GraspIt! (5.0), Newton Dynamics (5.6), and Flightmare (6.1) rank
lowest, consistent with their inactive, migrated, or low-activity status. Category
means are tightly bunched (7.3--7.8), so the rankings separate the top and
bottom tiers more clearly than the middle. An elicited-weight AHP ranking remains future work; an illustrative
equal-weight AHP ranking reproduces the same leaders.

\textbf{RQ3 (repository artifacts).} Practical artifacts are near-universal in the MPS sample
(Chapter~\ref{ch:comparison-research-software}):
installation instructions, installation automation, tutorials, user manuals,
API documentation, issue tracking, and version control are present for all 28
packages, which compares favourably with the medical imaging and LBM samples.
The recurring gaps are formal artifacts: documented user characteristics (1 of
28) and explicit coding standards (14 of 28).

\textbf{RQ4 (tool usage).} The measured MPS projects are tightly connected to
modern open-source tooling: git-based version control and issue tracking are
universal, continuous-integration evidence appears for 26 of 28 packages, and
documentation-generation tools are recorded for 12. The repository-based
inclusion rule partly explains this uniformity, but within the measured set tool
adoption is stronger than in the LBM comparison point and at least comparable
with medical imaging.

\textbf{RQ5 (principles, processes, and methodologies).} Process evidence is
visible but artifact-bound: 25 of 28 packages expose a defined development
process, 27 publish release notes, and 26 provide contribution guidance. What
repositories do not reveal---design rationale, requirements assumptions, and
internal decision making---remains under-documented in MPS software, the same
broad pattern reported for other research software domains.

\textbf{RQ6--RQ9 (pain points and practices).} These require
developer-interview evidence; Chapter~\ref{ch:developer-interviews}
records the planned evidence base, and
Chapter~\ref{ch:recommendations} gives artifact-based recommendations.
% TODO (after interviews): replace this paragraph with summarized answers to
% RQ6--RQ9 based on the approved interview records.

\section{Key Findings}
\label{sec:key-findings}

Four findings stand out from the artifact evidence. First, the practical
engineering infrastructure of MPS software is strong: every measured package
can be installed from documented instructions with some form of automation, and
only two installations broke during measurement, both recoverably. Second, the
weaknesses are concentrated in formal and explanatory artifacts rather than in
tooling: user assumptions, coding standards, and design rationale are the
recurring gaps. Third, community popularity and measured quality are different
signals: Gazebo leads on quality yet sits mid-table on stars,
and several highly starred packages place mid-table on quality,
so popularity indicators are not quality proxies. Fourth,
activity status only loosely predicts measured quality: packages that go
inactive decay at different rates depending on the engineering capital they
accumulated while active, as the contrast between OpenRAVE and GraspIt!
illustrates.

\section{Future Work}
\label{sec:future-work}

The most immediate direction is the developer-interview component: completing
recruitment, the interviews, and the analysis would answer RQ6--RQ9 and connect the measured artifact gaps to the
pain points developers actually experience.
A second direction is an elicited-weight AHP ranking, replacing the equal weights
of Chapter~\ref{ch:measurement-results} with weights obtained by pairwise
comparison of the qualities. Longitudinal repeat measurements would track how the
domain's practices change over time, and the same methodology could be extended
to adjacent robotics software domains.

One further direction concerns the measurement process itself, since completing
the template by hand is the most time-consuming part of the work. An exploratory
trial in which an AI coding agent installed, verified, and measured two physics
engines (Appendix~\ref{appendix:ai-assisted-measurement}) produced factual
entries tied to public sources but subjective scores that still needed human
review. With a person checking the evidence, this kind of agent-assisted
measurement could make it practical to cover more packages, or to repeat the
measurements over time, without straining the time budget that constrains a
manual study.
